\documentclass{book}

\usepackage[utf8]{inputenc}  % para que funcionen las tildes
\usepackage{amsmath}
\usepackage{amssymb}
\usepackage{amsthm}
\usepackage{letltxmacro}
\LetLtxMacro\amsproof\proof
\LetLtxMacro\amsendproof\endproof
\usepackage{txfonts}
\usepackage{lmodern}
\usepackage[dvipsnames]{xcolor}
\usepackage{thmtools, thm-restate}
\usepackage{datetime} % para la hora de compilación
\usepackage{graphicx}
\graphicspath{{images/}}
\usepackage[spanish,es-noquoting]{babel} % es-noquoting es para que funcione tikz
\usepackage{mathabx} % para \divides
\usepackage{centernot} % para \centernot\inculdes
\usepackage{wrapfig}
\usepackage{multicol}
\usepackage{subcaption}
\usepackage{xfrac}
\usepackage{glossaries}
\usepackage{standalone}
\usepackage{pgfmath}
\usepackage{pgfplots}
\usepackage{pdfpages}
\usepackage{tikz}
\usepackage{mathtools}
\usetikzlibrary{arrows.meta}
\usetikzlibrary{shapes}
\usetikzlibrary{decorations.text}
\usetikzlibrary{positioning}
\usetikzlibrary{external}
\tikzexternalize[prefix=./tikzbuild/]
\tikzexternalize % activate!

\AtBeginDocument{%
  \LetLtxMacro\proof\amsproof
  \LetLtxMacro\endproof\amsendproof
}

\definecolor{coolblack}{rgb}{0.0, 0.18, 0.39}
%\usepackage[a6paper,margin=5mm]{geometry}
\usepackage[a4paper, top=2.5cm,bottom=2.5cm,left=2.5cm,right=2.5cm]{geometry}

%\setlength{\parindent}{0pt}
\usepackage{parskip}

\usepackage{hyperref}

% ESTILOS DE DEFINICIONES Y TEOREMAS
\declaretheoremstyle[
	bodyfont=\normalfont,
	shaded={
		margin=8pt,
		bgcolor=White,
		rulecolor=Black,
		rulewidth=1pt
}]{mythm}

\declaretheoremstyle[
	bodyfont=\normalfont,
	shaded={
		margin=1em,
		bgcolor={rgb}{0.9,0.9,0.9}
}]{mydfn}

\declaretheoremstyle[
bodyfont=\normalfont,
spacebelow=1em,
spaceabove=1em,
]{myej}

\declaretheoremstyle[
bodyfont=\normalfont,
spacebelow=1em,
spaceabove=1em,
shaded={
    margin=8pt,
    bgcolor=White,
    rulecolor=coolblack,
    rulewidth=1pt
}]{thej}

\declaretheoremstyle[
bodyfont=\normalfont,
spacebelow=1em,
spaceabove=1em,
thmbox=M
]{myeg}

% DEFINICIONES DE ENTORNOS DE TEOREMAS
\declaretheorem[
	name=Teorema,
	refname={teorema,teoremas},
	Refname={Teorema,Teoremas},
	style=mythm
]{thm}

\declaretheorem[
	name=Corolario,
	refname={corolario,corolarios},
	Refname={Corolario,Corolarios},
	style=myej
]{cor}

\declaretheorem[
	name=Proposici\'{o}n,
	refname={proposici\'{o}n,proposiciones},
	Refname={Proposici\'{o}n, Proposiciones},
	sharenumber=thm,
	style=myej
]{pro}

\declaretheorem[
name=Lema,
refname={lema,lemas},
Refname={Lema, Lemas},
sharenumber=thm,
style=myej
]{lm}

%\newtheorem*{cor}{Corolario}
\newtheorem*{lem}{Lema}

\declaretheorem[
	name=Definici\'{o}n,
	refname={definici\'{o}n,definiciones},
	Refname={Definici\'{o}n,Definiciones},
	style=mydfn
]{dfn}

\declaretheorem[
name=Ejercicio,
refname={ejercicio,ejercicios},
Refname={Ejercicio,Ejercicios},
style=myej,
numbered=no
]{ex}

\declaretheorem[
name=Ejercicio propuesto,
refname={ejercicio,ejercicios},
Refname={Ejercicio,Ejercicios},
style=thej,
]{th_ex}

\declaretheorem[
name=Ejemplo,
refname={ejemplo,ejemplos},
Refname={Ejemplo,Ejemplos},
style=myeg
]{eg}

\declaretheorem[
name=Observaci\'{o}n,
refname={observaci\'{o}n,observaciones},
Refname={Observaci\'{o}n,Observaciones},
style=myej,
numbered=no
]{obs}

\declaretheorem[
name=Notaci\'{o}n,
refname={notaci\'{o}n},
Refname={Notaci\'{o}n},
style=myej,
numbered=no
]{nota}

\makeglossaries
\newglossaryentry{di}{name={DI},description={Dominio de integridad (ver \autoref{dfn:dominiointegridad})}}
\newglossaryentry{dip}{name={DIP},description={Dominio de ideales principales (ver \autoref{dfn:dominioidealesprincipales})}}
\newglossaryentry{dfu}{name={DFU},description={Dominio de factorización única (ver \autoref{dfn:dfu})}}
\newglossaryentry{de}{name={DE},description={Dominio euclídeo (ver \autoref{dfn:de})}}


\setcounter{chapter}{0}
%COMANDOS ÚTILES PARA DERIVADAS
\AtBeginDocument{\renewcommand{\d}{\mathrm{d}}}
\newcommand{\Dd}[1]{\frac{\mathrm{d}}{\mathrm{d}#1}}
\newcommand{\dd}[2]{\frac{\mathrm{d}#1}{\mathrm{d}#2}}
\newcommand{\mbf}[1]{\mathbf{#1}}
\newcommand{\mtc}[1]{\mathcal{#1}}
% COMANDOS ÚTILES PARA LA TEORÍA DE GRUPOS
%\newcommand{\normsub}{\mathbin{\triangleleft}}
\newcommand{\normsub}{\lhd}
\newcommand{\uds}[1]{\mathcal{U}(#1)}
\newcommand{\inv}[1]{#1^{-1}}
\newcommand{\ima}{\text{Im }}
\newcommand{\isom}{\simeq}
\newcommand{\autom}[1]{\text{Aut}(#1)}
\newcommand{\gen}[1]{\langle#1\rangle}
\newcommand{\biy}[1]{\text{Biy}(#1)}

%\DeclareMathSymbol{\varprod}{\mathop}{largesymbolsA}{16}

\newcommand{\hr}{\rule{\textwidth}{.4pt}}
\newcommand{\N}{\mathbb{N}}
\newcommand{\Z}{\mathbb{Z}}
\newcommand{\R}{\mathbb{R}}
\newcommand{\Q}{\mathbb{Q}}
\newcommand{\C}{\mathbb{C}}
\newcommand{\A}{\mathbb{A}}
\newcommand{\F}{\mathbb{F}}
\newcommand{\J}{\mathbb{J}}
\newcommand{\B}{\mathbb{B}}
\DeclareMathOperator{\diag}{diag}
\DeclareMathOperator{\traza}{traza}
\renewcommand{\P}{\mathcal{P}}
\renewcommand{\S}{\mathbb{S}}
\newcommand{\ZnZ}{\mathbb{Z}/n\mathbb{Z}}
\newcommand{\ZmZ}{\mathbb{Z}/m\mathbb{Z}}
\newcommand{\0}{\mathbf{0}}
\newcommand{\1}{\mathbf{1}}
\newcommand{\gr}{\text{gr }}
\newcommand\ninf[1]{\left|\left|#1\right|\right|_\infty}

%%%% MATRICES
\newcommand\mx[1]{\begin{bmatrix*}#1\end{bmatrix*}}

\renewcommand\qedsymbol{$\diamondsuit$}

\title{Apuntes de L\'ogica M\'atematica}
\author{Rafael S\'{a}nchez S\'{a}nchez}
\begin{document}
\begin{titlepage}
	\includepdf{portada.pdf}
\end{titlepage}


Revisión del \today $ $ a las \currenttime.

\begin{center}

\end{center}

\tableofcontents
\part{Sintaxis y sem\'antica de primer orden}
% !TeX root = ../logica-matematica.tex

\chapter{Estructuras y lenguajes de primer orden}
\section{Introducción}
\begin{dfn}[Estructura]\label{dfn:estructura}
    Una \textbf{estructura} (de primer orden) $\mtc{A}$ consta de un conjunto no vacío $A$ (universo) y un conjunto de funciones, relaciones y elementos del universo.
\end{dfn}
Puede parecer una definición algo abstracta, así que vamos a ver algunos ejemplos:
\begin{eg}[Estructuras. Ejemplos]
    \begin{itemize}
        \item $ \mtc{G} = \gen{G, \cdot, 1} $
        \item $ R = \gen{\R,\  +,\  -,\  \cdot,\  0,\  1,\  \leq} $
        \item $ R_* = \gen{\R,\  +,\  -,\  \cdot,\  0,\  1,\ \left. \sin \right|_{[0, 1]},\  \exp}$
    \end{itemize}
\end{eg}

\begin{dfn}[Lenguaje]\label{dfn:lenguaje}
    Un \textbf{lenguaje} (de primer orden) consta de:\\
    \begin{itemize}
        \item Para cada $n \in \N^*$ un conjunto $\mtc{F}_n$ de símbolos de funciones n-arias.
        \item Para cada $m \in \N^*$ un conjunto $\mtc{R}_m$ de símbolos de relaciones m-arias.
        \item Un conjunto $\mtc{C}$ de constantes.
    \end{itemize}
    Además de conjuntos de símbolos lógicos, variables $\{v_i\}_{i\in \N}$, cuantificadores $(\exists,\ \forall)$, conectores $(\land, \lor, \iff, \implies)$, símbolos de igualdad $=$, y paréntesis $()$. No se suelen especificar en la declaración del lenguaje.
\end{dfn}

\begin{eg}[Lenguajes. Ejemplos]\label{eg:lenguajes}
    \begin{itemize}
        \item $L = \left\{\ast,\ e\right\}$
        \item $L = \left\{\cdot,\ 1\right\}$
        \item $L = \left\{+,\ 0\right\}$
        \item $L = \left\{\cdot,\ ^{-1},\ 1\right\}$
    \end{itemize}
    Además, para hablar del comportamiento de los reales vamos a usar el lenguaje:
    $$
        L = \left\{+,\ -,\ \cdot,\ 0,\ 1,\ \leq \right\}
    $$
    Es importante destacar que cuando escribimos $+$ en la declaración del lenguaje no nos referimos a la función: $+: \R^2 \to \R$ si no a un símbolo que luego interpretamos como tal función en la estructura.
\end{eg}

\section{L-estructuras}
\begin{dfn}[L-estructura]
    Dado un lenguaje \textsuperscript{\ref{dfn:lenguaje}} $L$, una \textbf{L-estructura} $\mtc{A}$ o una \textbf{interpretaci\'on de L} consta de:
    \begin{itemize}
        \item un universo $A \neq \varnothing$
        \item una función n-aria $f^\mtc{A} : A^n \to A$ para cada símbolo de función $f \in \mtc{F}_n$
        \item una relación m-aria $R^\mtc{A} \subseteq A^m$ para cada símbolo de relación $R \in \mtc{R}_m$
        \item un elemento $e^\mtc{A}\in A$ para cada constante $c \in \mtc{C}$.
    \end{itemize}
\end{dfn}

\begin{nota}
    En la definición anterior:\\
    $L = \mtc{F} \cup \mtc{R} \cup \mtc{C}$ (lenguaje), $\mtc{F}=\bigcup_{n\in \N^*} \mtc{F}_n$ (símbolos de función), $\mtc{R}=\bigcup_{m\in \N^*} \mtc{R}_m$ (símbolos de relación), $\mtc{C}$ constantes.\\
    Además, para cada símbolo $s\in L$, $s^\mtc{A}$ es la interpretación de $s$ a $\mtc{A}$.
\end{nota}

\begin{eg}[L-estructura. Ejemplos]
    \begin{itemize}
        \item $R = \gen{\R,\ +^R,\ -^R,\ \cdot^R,\ \leq^R,\ 0^R,\ 1^R}$ (interpretación del lenguaje de los reales con el universo $\R$).
        \item $\mtc{A} = \gen{A,\ +^\mtc{A},\ -^\mtc{A},\ \cdot^\mtc{A},\ \leq^\mtc{A},\ 0^\mtc{A},\ 1^\mtc{A}}$.
        $$
            +,\ -,\ \cdot \in \mtc{F}_2 \implies +^\mtc{A},\ -^\mtc{A},\ \cdot^\mtc{A}: A^2 \to A
        $$
        $$
            \leq \in R_2 \implies \leq^\mtc{A} \in A^2;\ \ \ \ \ 0, 1 \in \mtc{C} \implies 0^\mtc{A}, 1^{\mtc{A}} \in A
        $$
    \end{itemize}
\end{eg}

\begin{obs}
    Podríamos intentar interpretar el lenguaje de los reales del ejemplo \ref{eg:lenguajes} con el universo $A = \C$, sin embargo, aunque podemos interpretar $+, -, \cdot, 0$ y $1$ de la forma habitual, no existe una interpretación de $\leq$ en $\C$.
\end{obs}

\begin{eg}[Lenguajes comunes]
    \begin{itemize}
        \item $L_\varnothing = \{\}$. Es el \textit{lenguaje vacío}, sigue teniendo símbolos generales. Sirve para expresar propiedades tales como: \textit{Existen tres elementos} ($\exists x_1, x_2, x_3$).
        \item $L_< = \{<\}$. Lenguaje para \textit{conjuntos ordenados}. Con $<\ \in R_2$.
        \item $L_{grupos} = \{+, -, 0\}$ (aditivo), $\{\cdot, ^{-1}, 1\}$ (multiplicativo). Lenguaje para \textit{grupos}. Con $+,\cdot\ \in \mtc{F}_2$; $-, ^{-1}\ \in \mtc{F}_1$; $0, 1 \in \mtc{C}$.
        \item $L_{cuerpos} = \{+, -, \cdot, 0, 1\}$. Lenguaje para \textit{cuerpos}.
        \item $L_{aritm\'etica} = \{+, \cdot, 0, 1, \leq\}$. Lenguaje para la \textit{aritmética}. También podemos considerar añadir otro símbolo de función $S$, cuya interpretación natural sería la función \textit{sucesor}.
        \item $L_{conj} = \{\in\}$. Lenguaje para \textit{conjuntos}. Todo se puede escribir con este lenguaje.
    \end{itemize}
\end{eg}

\subsection{Subestructuras}
\begin{dfn}[Subestructura de una L-estructura]
    Sean $\mtc{B}, \mtc{A}$ L-estructuras (con universos $B$ y $A$ respectivamente), decimos que $\mtc{A}$ es una subestructura de $\mtc{B}$ ($\mtc{A} \subseteq \mtc{B}$) si:
    \begin{itemize}
        \item $f^\mtc{A} = \left. f^\mtc{B} \right|_{A^n}$ para cada $f \in \mtc{F}_n$.
        \item $R^\mtc{A} = R^\mtc{B} \cap A^m$ para cada $R \in R_m$.
        \item $C^\mtc{A} = C^\mtc{B}$
        \item $A \subseteq B$
    \end{itemize}
\end{dfn}
\begin{eg}[Subestructuras. Ejemplos]
    Sean los lenguajes: $L_1 = \{+, 0\}$, $L_2 = \{+, -, 0\}$. Vamos a considerar las $L_1$-estructuras:
    $$
        W = \gen{\N, +^W, 0^W};\ Z = \gen{\Z, +^Z, 0^Z}
    $$
    Donde es fácil ver que se cumplen las condiciones de subestructura y podemos afirmar que $W \subseteq Z$. Sin embargo, si consideramos las $L_2$-estructuras:
    $$
        W' = \gen{\N, +^{W'}, -^{W'}, 0};\ Z' = \gen{\Z, +^{Z'}, -^{Z'}, 0}
    $$
    Vamos a dar una definición de $-^{W'}$ ya que el opuesto no está bien definido en $\N$. $-^{W'}: n \to 0$. Con esta interpretación es fácil ver que no se cumple que $-^{W'} = \left. f^{Z'} \right|_{\N}$ ya que:
    $$
        -^{W'}(2) = 0 \text{ y sin embargo } -^{Z'}(2)=-2
    $$
\end{eg}
\begin{obs}
    Consideremos el lenguaje $L_< = \{<\}$ de conjuntos ordenados, y las L-estructuras $\mtc{A} = \gen{A, <}$, $\mtc{B} = \gen{B, <}$ (que sólo tienen una relación). Es fácil ver que:
    $$
        \mtc{A} \subseteq \mtc{B} \iff A \subseteq B \text{ (suponemos $<^\mtc{A} = <^\mtc{B} \cap\ A^2$)}
    $$
\end{obs}

\section{Homomorfismos}
\begin{dfn}[Homomorfismo y monomorfismo]
    Sean $\mtc{A}$, $\mtc{B}$ L-estructuras y $\varphi: A \to B$ una función, decimos que $\varphi: \mtc{A} \to \mtc{B}$ es un \textbf{homomorfismo} de $\mtc{A}$ en $\mtc{B}$ si:
    \begin{enumerate}
        \item $\forall a_1, \ldots, a_n \in A$ se cumple: $\varphi(f^\mtc{A}(a_1, \ldots, a_n)) = f^\mtc{B}(\varphi(a_1), \ldots, \varphi(a_n))$
        \item $\forall a_1, \ldots, a_n \in A$ se cumple: $(a_1, \ldots, a_m) \in R^\mtc{A} \implies (\varphi(a_1), \ldots, \varphi(a_m) \in \mtc{B})$
        \item $\forall c \in \mtc{C}$ se cumple: $\varphi(c^\mtc{A}) = c^\mtc{B}$
    \end{enumerate}
    Además, si la segunda condición resulta ser necesaria y suficiente, $\varphi$ es un \textbf{monomorfismo}.
\end{dfn}
\begin{eg}[Ejemplo de monomorfismo]
    Sean $\mtc{A} = \gen{A, +, 0}$, $\mtc{B} = \gen{B, +, 0}$ y $\varphi(a+b) = \varphi(a) + \varphi(b);\ \varphi(0)=0$. Es fácil ver que $\varphi$ es un homomorfismo. Además, la única relación es la igualdad (que existe aunque no se especifique) y como $\varphi(a) = \varphi(b) \iff a = b$ estamos ante un monomorfismo.
\end{eg}

\part{Ap\'{e}ndices}

%\renewcommand{\chaptername}{Apéndice}
%\renewcommand{\thechapter}{\Alph{chapter}}
%\setcounter{chapter}{0}

%% !TeX root = ../apuntes-ea.tex

\chapter{Ejercicios}

\section{Hoja 1}

\begin{ex}[H1.2]
	Sean $a,b,c \in G = (-1,1)$. Probamos las propiedades de los grupos.
	\begin{itemize}
		\item \textbf{Asociatividad:}
		\begin{align*}
			(a \ast b) \ast c = \left(\frac{a+b}{1+ab}\right) \ast c = \frac{\left(\frac{a+b}{1+ab}\right) + c}{1 + \left(\frac{a+b}{1+ab}\right)c} = \frac{\frac{a+b+c+abc}{1+ab}}{\frac{1+ab+ac+bc}{1+ab}} = \frac{a+b+c+abc}{1+ab+ac+bc} \\
			a \ast (b \ast c) = a \ast \left(\frac{b+c}{1+bc}\right) = \frac{a+\left(\frac{b+c}{1+bc}\right)}{1+a\left(\frac{b+c}{1+bc}\right)} = \frac{\frac{a + b + c + abc}{1+bc}}{\frac{1+ab+ac+bc}{1+bc}} = \frac{a+b+c+abc}{1+ab+ac+bc}
		\end{align*}
		\item \textbf{Elemento neutro:} es el $0$ ya que $x\ast 0 = \frac{x+0}{1+x\cdot 0} = \frac{x}{1} = x$ y además $0 \ast x = \frac{0+x}{1+0\cdot x} = \frac{x}{1} = x$
		\item \textbf{Elemento inverso:} la ecuación
		\begin{align*}
			x \ast \inv{x} = 0 \iff \frac{x + \inv{x}}{1+x\inv{x}} = 0 \iff \inv{x} = -x
		\end{align*}
		siempre tiene solución y ocurre lo mismo para la ecuación $\inv{x} \ast x = 0 \iff \inv{x} = -x$
		\item \textbf{Clausura:} tenemos que probar que si $x,y \in (-1,1)$ entonces $x \ast y \in (-1,1)$. Consideramos $f(x,y) = x \ast y = \frac{x+y}{1+xy}$. Derivando tenemos que $\nabla f(x,y) = (\frac{1}{(1+xy)^2},\ \frac{1}{(1+xy)^2}) \neq 0, \forall x,y \in [-1,1]\times[-1,1]$. Si el máximo no se alcanza en ningún sitio de dentro del cuadrado $(-1,1)\times(-1,1)$ se tendrá que alcanzar en el borde.
		\begin{itemize}
			\item Fijado $x = 1$ tenemos que $f(1,y) = \frac{1+y}{1+y} = 1 \implies f(1, -1 < y < 1) < 1$ porque si $f(1, -1 < y < 1)$ tomara un valor mayor que $1$ habría un máximo en $(-1,1)\times(-1,1)$ y esto no puede ser pues $\nabla f$ no se anula en el cuadrado.
			\item Fijado $x = -1$ tenemos que $f(-1, y) = \frac{y-1}{1-y} = -1 \implies f(-1, -1 < y < 1) > -1$ por la misma razón que antes.
			\item Hacemos lo mismo fijando la $y$ y variando la $x$.
		\end{itemize}
		En el borde (que no está incluido) se alcanzan máximo y mínimo que acotan a $f$ en el cuadrado:
		\begin{align*}
			-1 < f(x,y) = x \ast y < 1,\quad \forall x, y \in G
		\end{align*}
	\end{itemize}
\end{ex}

\begin{ex}[H1.3]Hallar los inversos de los siguientes elementos, cada uno en su grupo correspondiente:
	\begin{enumerate}
		\item $o(\overline{11})$ en $\uds{\Z^\ast/23\Z}$ es $22$ porque $23 \cdot 22 \equiv 1 \mod 23$
		\item $o(\overline{5})$ en $\uds{\Z^\ast/31\Z}$ es $3$ porque $5 \cdot 3 \equiv 1 \mod 31$
	\end{enumerate}
\end{ex}

\begin{ex}[H1.33]
	\label{ex:h1.33}
	Sea $G$ un grupo. Suponed que existe un único $a \in G$ de orden 2. Demostrad que $a \in Z(G)$.
\end{ex}

\begin{proof}
	Recordamos que $a \in Z(G) \iff ga = ag,\ \forall g \in G$. Definimos el isomorfismo de conjugación $\phi_g (x) = gx\inv{g}$ para algún $g$. Como $\phi_g$ es isomorfismo lleva elementos de orden $n$ en elementos de orden $n$. Entonces $\phi_g(a) = a$ ya que $a$ es el único elemento de orden 2. Por tanto $ga\inv{g} = a \implies ga = ag \implies a \in Z(G)$.
\end{proof}

\section{Hoja 2}

\begin{ex}[H2.1]
	\label{ex:h2.1}
	Se considera el tercer grupo diédrico $D_3$. Se pide hallar lo siguiente:
	\begin{enumerate}
		\item Las clases de conjugación de cada uno de sus elementos.
		\begin{proof}
			Las clases dan una partición del grupo. Si un elemento pertenece a una clase, entonces la clase de ese elemento también es la clase a la que pertenece.
			\begin{itemize}
				\item $cl(e) = \{e\}$
				\item $cl(B)$? Sabemos que $|cl(B)| = [G:C(B)]$. Sabemos que $\gen{B} =\{1, B, B^2\} \subset C(B)$ luego $|C(B)| \geq 3$. Si hubiera más elementos en $C(B)$ tendríamos que $|C(B)| = 6$ pues $C(B) < D_3$. Esto no ocurre porque sabemos que $B$ no conmuta con todos los demás elementos. Por ejemplo $BA \neq AB$. Por tanto $|C(B)| = 3 \implies |cl(B)| = [D_3:C(B)]= 6/3 = 2$. Es claro que $B \in cl(B)$. Además, como $cl(B)$ contiene elementos transformados por el isomorfismo conjugación sabemos que el otro elemento que hay tiene orden 3. El único elemento que queda de orden 3 es $B^2 \implies cl(B) = \{B, B^2\}$.
				
				\item $cl(A)?$ Sabemos que $A$ no conmuta con todos ($A \not\in Z(D_3)$) luego $|C(A)| < 6$. Sabemos que $\gen{A} = \{1, A\} < C(A)$. Además, como $C(A)$ es un (sub)grupo sabemos que no puede haber más elementos porque si los hubiera, $|\gen{A}| \divides |C(A)| \implies C(A) \geq 6$ pero ya hemos visto que no puede ser. Es decir que $|cl(A)| = [D_3:D(A)] = 6 / 2 = 3$. Por tanto $cl(A)$ incluye los 3 elementos que nos quedan: $cl(A) = \{A, AB, AB^2\}$.
			\end{itemize}
		\end{proof}
	
		\item Los elementos de $\text{Int}(D_3)$.
		\item Los centralizadores $C_{D_3}(x)$ para cada $x \in D_3$
		\item Los normalizadores $N(H)$ para cada $H < D_3$.
	\end{enumerate}
\end{ex}

\begin{ex}[H2.2]
	
	\begin{proof}
		Obtenidas las clases en el ejercicio \nameref{ex:h2.1} se verifica que $|D_3| = |cl(e)| + |cl(B)| + |cl(A)| = 1 + 2 + 3 = 6$
	\end{proof}
	
\end{ex}

\begin{ex}[H2.6]
	Sea $G$ un grupo. ¿Verdadero o falso?
	\begin{enumerate}
		\item $H < G$ y $H$ conmutativo implica $H \normsub G$.
		\item $H < G$ y $|H| = 2$ implica $H \normsub G$.
		\item Si $\varphi: G \to G_1$ es un homomorfismo de grupos, entonces $\ima \varphi \normsub G$
		\item Si $H \normsub K$ y $K \normsub G$ entonces $H \normsub G$
		\item Si $H \normsub G$ y $|H| = m$ entonces $H$ es el único subgrupo de $G$ de orden $m$.
		\item Si $H \normsub G$ entonces $H < Z(G)$.
		\begin{proof}[FALSO]
			Contraejemplo: En $G = D_4$ tomamos $H = \gen{B^2} = \{1, B, B^2, B^3\} \not\subset Z(D_4) = \{1, B^2\}$.
		\end{proof}
	\end{enumerate}
\end{ex}

\begin{ex}[H2.10]
	\begin{proof}
		Fijado $n$ y definida $\alpha_n : G \to G,\ x \mapsto x^n$ tenemos que $\alpha_n$ es un homomorfismo de grupos. Además podemos expresar $H_2 = \ker \alpha_n \implies H_2 \normsub G$. Además también tenemos que $H_1 = \ima \alpha_n < G$. Veamos que $H_1 \normsub G$. Es decir, que $gH_1\inv{g} = H_1,\ \forall g \in G$. Para ello tomamos $x_1^n \in H_1$ y lo conjugamos $gx_1^n\inv{g} = (gx_1\inv{g})^n$ por ser $\alpha$ homomorfismo de grupos. En particular $(gx_1\inv{g})^n \in \ima \alpha \implies (gx_1\inv{g})^n \in H_1 \implies (gx_1\inv{g})^n = x_2^n$ para algún $x_2 \in H_1 \implies H_1 \normsub G$.
	\end{proof}
\end{ex}

\begin{ex}[H2.13] Si $A$ es un grupo abeliano con $n$ elementos y $k$ es un entero primo con $n$, demostrad que la aplicación $\varphi : A \to A$ definida por $\varphi(a) = a^k$ es un isomorfismo.
	\begin{itemize}
		\item $\varphi$ homomorfismo de grupos.
		\begin{proof}
			\begin{align*}
				\varphi(a)\varphi(b) = a^kb^k = (ab)^k = \varphi(ab)
			\end{align*}
		\end{proof}
		\item $\varphi$ biyectiva $\iff \varphi$ inyectiva ya que dominio y codominio coinciden
		\begin{proof}
			$\ker \varphi = \{a \in A \mid \varphi(a) = a^k = 1\}$. Probaremos que $a^k = 1 \iff a = 1$ y por tanto que $\ker \varphi = \{1\} \implies \varphi$ inyectiva. Sabemos que $a^k = 1 \iff o(a^k) = 1$. Sea $t = o(a) \divides n$. Distinguimos dos casos
			\begin{itemize}
				\item Si $t = 1$ entonces $a = 1$ y ya está
				\item Si $t > 1$ entonces $o(a^k) = \frac{t}{mcd(k, t)} = \frac{t}{1} > 1$ contradicción. Luego necesariamente $t = o(a) = 1$.
			\end{itemize}
		\end{proof}
	\end{itemize}
\end{ex}

\begin{ex}[H2.22]
	\label{ex:h2.22}
	Demostrad que si $G$ es un grupo no conmutativo y tiene orden $p^3$ ($p$ un número primo) entonces $Z(G)$ tiene orden $p$.
	
	\begin{proof}
		Sabemos que $Z(G) < G \implies |Z(G)| \divides |G| \implies |Z(G)| \in \{1, p, p^2, p^3\}$
		\begin{itemize}
			\item $|Z(G)| \neq p^3$ porque en tal caso $G$ sería conmutativo
			\item $|Z(G)| \neq 1$ porque $G$ es un p-grupo y por tanto su centro no es el trivial.
			\item Si $|Z(G)| = p^2$ entonces $|G/Z(G)| = p \implies G/Z(G)$ es cíclico lo que no es posible si $G$ no es abeliano.
		\end{itemize}
	
		Por descarte concluimos que $|Z(G)| = p$.
	\end{proof}
\end{ex}


\begin{ex}[H2.25]
	Sabemos que $\autom{\Z/12\Z} \isom \uds{\Z^\star/12\Z}$ donde $\Z^\star/12\Z$ es el grupo multiplicativo $(\{\overline{1}, \overline{2}, \overline{3}, \dots,  \overline{11}\}, \cdot)$. Queda $\uds{\Z^\star/12\Z} = (\{\overline{1}, \overline{5}, \overline{7}, \overline{11}\}, \cdot )$ y además da la casualidad que $\forall x \in \uds{\Z^\star/12\Z},\ o(x) = 2$ (todos los elementos son su propio inverso) por lo que no tenemos restricciones al definir $f: \Z/2\Z \to \autom{\Z/12\Z}$:
	\begin{align*}
		f:\Z/2\Z & \to \autom{\Z/12\Z} \isom \uds{\Z^\star/12\Z} \\
		e = \overline{0} &\mapsto 1 \\
		\overline{1} &\mapsto \{\overline{1}, \overline{5}, \overline{7}, \overline{11}\}
	\end{align*}
\end{ex}

\begin{ex}[H2.26]
	Sea $|G_1| = m, |G_2| = n,\ mcd(m, n) = 1$. Si $f:G_1 \to G_2$ es h. de g. sabemos que $o(f(a)) \divides o(a),\ \forall a \in G_1$. Además $o(a) \divides m \land o(f(a)) \divides n$ por el teorema de Lagrange (\ref{thm:lagrange}).
	\begin{align*}
		\begin{cases}
		o(a) \divides m \land o(f(a)) \divides n \\
		mcd(m, n) = 1 \\
		o(f(a)) \divides o(a)
		\end{cases} \implies o(a) = o(f(a)) = 1, \forall a \in G_1
	\end{align*}
	Por lo que solo puede haber un homomorfismo entre ellos y además es el trivial $f(a) = e_{G_2}$.
\end{ex}

\begin{ex}[H2.19]
	Definimos una función $f:[0, 2\pi] \subset \R \to \mathbb{S}^1, \alpha\mapsto \cos \alpha + i \sin \alpha$. Esta función tiene la propiedad de que $f(\alpha) \cdot f(\alpha') = \cos(\alpha + \alpha') + i\sin(\alpha + \alpha') = f(\alpha + \alpha')$ y por tanto es un h. de g.\footnote{A la izquierda (en $\R$) sumamos pero a la derecha (en $\mathbb{S}^1$) multiplicamos.} entre $R$ y $\mathbb{S}^1$.
	
	Un elemento de $\cos \alpha + i \sin \alpha \in \mathbb{S}^1$ es de torsión $\iff \exists n \mid (\cos \alpha + i \sin \alpha )^n = 1$. Ahora bien $(\cos \alpha + i \sin \alpha )^n = \cos n\alpha + i \sin n\alpha = 1 \iff n\alpha = k 2\pi$. 
\end{ex}

\section{Hoja 4}

\begin{ex}[H4.11] Hallar los subgrupos de Sylow de $S_5$.
	Sabemos que $|S_5| = 5! = 2^3\cdot3\cdot 5$ y por \nameref{thm:sylow1} tenemos lo siguiente:
	\begin{itemize}
		\item $\exists P_2,\ |P_2| = 2^3 = 8$.
		\item $\exists P_3,\ |P_3| = 3$. Además en $S_5$ hay $\binom{5}{3}2! =20 $ 3-ciclos y en cada $gP_3\inv{g} = \{1, a, a^2 \mid o(a) = 3\}$ hay 2 elementos de orden 3 distintos. Además, $g_1P_3\inv{g_1} \cap g_2P_3\inv{g_2} = \{e\}$ porque si su intersección fuera más grande entonces serían el mismo subgrupo (porque son cíclicos). Es por esto que tenemos que repartir 40 elementos dando 2 a cada 3-grupo con lo que obtenemos $n_3 = 20 / 2 = 10$ 3-subgrupos de Sylow en $S_5$.
		\item $\exists P_5,\ |P_5| = 5$. Además en $S_5$ hay $4! = 24$ 5-ciclos (elementos de orden 5) y en cada $gP_5\inv{g} = \{1, a, a^2, a^3, a^4 \mid o(a) = 5\}$ tenemos $4$ elementos de orden $5$ distintos. Además, $g_1P_5\inv{g_1} \cap g_2P_5\inv{g_2} = \{e\}$ porque si su intersección fuera más grande entonces serían el mismo. Así, tenemos $24$ 5-ciclos a repartir entre los diferentes $gP_5\inv{g}$ dando 4 5-ciclos a cada 1. Por tanto tenemos $n_5 = 24/4 = 6$ 5-subgrupos de Sylow en $S_5$.
	\end{itemize}
\end{ex}

\begin{ex}[H4.18]
	Demostrar que todo grupo de orden $|G| = 5^3 \cdot 7^3$ tiene un subgrupo normal de orden 125.
	
	\begin{proof}
		\nameref{thm:sylow1} $\implies \exists P_5 < G,\ |P_5| = 125$. \nameref{thm:sylow3} $\implies n_5 \divides 7^3 \land n_5 \equiv 1 \mod 5$ es decir $n_5 \in \{1, 7, 49, 343\} \land n_5 \in \{1, 6, 11, \dots\}$. Como ni $49$ ni $343$ son conguentes con $1$ módulo 5 tenemos que $n_5 = 1 \implies P_5 \normsub G$.
	\end{proof}
\end{ex}

\begin{ex}[H4.20] Hallad todos los grupos abelianos de órdenes 36, 64, 96 y 100.
	\begin{enumerate}
		\item $|G| = 36 = 2^2\cdot3^2$
		\begin{proof}
			\nameref{thm:sylow1} $\implies \exists P_2,\ |P_2| = 4 \land \exists P_3,\ |P_3| = 9$. Además $G$ abeliano $\implies P_2, P_3 \normsub G \implies G \isom P_2 \times P_3$. Estudiamos los grupos de orden 4 y de orden 9
			\begin{itemize}
				\item $|P_2| = 4$ entonces $P_2 \isom \Z/2\Z \times \Z/2\Z \lor P_2 \isom \Z/4\Z$
				\item $|P_3| = 9$ entonces $P_3 \isom \Z/3\Z \times \Z/3\Z \lor P_3 \isom \Z/9\Z$
			\end{itemize}
			Como $G \isom P_2 \times P_3$ tenemos 4 posibles grupos abelianos de orden 36.
		\end{proof}
	\end{enumerate}
\end{ex}

\begin{ex}[H4.22]
	Hallar todos los grupos abelianos de orden 175.
	
	\begin{proof}
		$|G| = 5^2\cdot 7$. Por el \nameref{thm:sylow1} tenemos que $\exists P_5, P_7 < G$ con $|P_5| = 25,\ |P_7| = 7$ y además por ser $G$ abeliano tenemos que $P_5, P_7 \normsub G \implies G \isom P_5 \times P_7$. Estudiamos los grupos de orden 25 y de orden 7:
		\begin{itemize}
			\item $|P_5| = 25 \land P_5$ abeliano $\implies P_5 \isom \Z/25\Z \lor P_5 \isom \Z/5\Z \times \Z/5\Z$. En ambos casos $P_5$ es producto directo de cíclicos pues $\Z/n\Z$ es cíclico.
			\item $|P_7| = 7 \land P_7$ abeliano $\implies P_7 \isom \Z/7\Z$. Ocurre lo mismo que con $P_5$.
		\end{itemize}
		Concluimos que $G \isom \Z/25\Z \times \Z/7\Z \lor G \isom \Z/5\Z \times \Z/5\Z \times \Z/7\Z$. Los dos casos son abelianos por ser producto directo de grupos cíclicos.
	\end{proof}
\end{ex}

\begin{ex}[H4.23]
	¿Cuántos elementos de orden 3 puede tener un grupo abeliano de orden 36?
	
	\begin{proof}
		$|G| = 36 = 2^2 3^2$. \nameref{thm:sylow1} $\implies \exists P_2, P_3 < G,\ |P_2| = 4,\ |P_3| = 9$. Estudiamos los grupos de órdenes 4 y 9:
		\begin{itemize}
			\item $|P_3| = 9 \implies P_3 \isom \Z/9\Z \lor P_3 \isom \Z/3\Z \times \Z/3\Z$
			% TODO: terminar
		\end{itemize}
	\end{proof}
\end{ex}

\chapter{\'{I}ndices}

%%TODO por qué no podemos utilizar only= para hacer whitelisting en lugar de blacklisting? que alguien lo arregle parfavar
\renewcommand{\listtheoremname}{Lista de definiciones}
\listoftheorems[ignore={thm,eg,pro,cor,obs,nota,lm,ex,th_ex}]

\renewcommand{\listtheoremname}{Lista de teoremas}
\listoftheorems[onlynamed,ignore={dfn,eg,cor,obs,nota,ex,th_ex}]

\renewcommand{\listtheoremname}{Lista de ejemplos}
\listoftheorems[onlynamed,ignore={dfn,thm,pro,cor,obs,nota,lm,ex,th_ex}]

\renewcommand{\listtheoremname}{Lista de ejercicios}
\listoftheorems[ignore={dfn,thm,pro,cor,obs,nota,lm,eg}]

\printglossaries

\bibliographystyle{alpha}

\end{document}
